\scanpage{154}
\[
\dim_S\underline G_H(q)=d^2+d\rho-d-\sum\frac{d_i(d_i-1)}2
\]

Par exemple
\[
\dim_S\underline G_{\{d\}}(q)=\frac{d(d+1)}{2}+d(\rho\deleted{\unsure{+}}-1)\qquad(\rho=N-2d)
\]
\deleted{\unsure{Ai|Ainsi}}Ainsi, si $\unsure{q|f}$ est de rg $N=2n$, et $d=\{n\}$, \inserted{(donc $\rho=0$)} on trouve
\[
\frac{\deleted{\unsure{d}}n(n+1)}{2}\deleted{\unsure{+|-}}-n=\frac{n(n-1)}{2}
\]

\begin{center}\rule{.14\linewidth}{.4pt}\end{center}

Supposons $N=2\underset{\deleted{\unsure{0}}}{n}+1$ --- de sorte que
\[
\mathrm O(q)=\mathrm{SO}(q)\times\mu_2,
\]
\deleted{On va vérifier} \inserted{où}
\[
\left\{\begin{aligned}
&\mu_2=\mathrm O(q)\cap G_m=\mathrm Z(\mathrm O(q))\\
&\underline{\mathrm{SO}}(q)=\underline{\mathrm O}(q)\cap\underline{\mathrm{Sl}}(E)
\end{aligned}\right.
\]
\deleted{\unsure{$\Rightarrow$ il résulte}} \deleted{que $\mathrm{SO}(q)$ est} $\mathrm O(q)$ opère sur les $\underline G_H(q)$ avec $\mu_2$ opérant trivialement, donc il opère via le quotient $\mathrm{SO}(q)$ --- donc les $\underline G_H$ sont aussi des espaces homogènes sous $\mathrm{SO}(q)$. \deleted{\unsure{Pour le|Pour}} \deleted{\unsure{dès que}} \inserted{Pour} $\forall\,D\in G_H$, \inserted{soit $\underline{\mathrm{SO}}_D(q)$} le stabilisateur dans $\underline{\mathrm{SO}}(q)$, \deleted{donc}. On a donc
\[
\underline{\mathrm{SO}}_D(q)=\underline{\mathrm O}_D(q)\cap\underline{\mathrm{Sl}}(E)
\]
Considérons l'hom canonique
\[
\underline{\mathrm O}_D(q)\xrightarrow{\alpha}\deleted{\prod_{1\leq i\leq\unsure{n|r|\nu}}\underline{\mathrm{Gl}}(\unsure{R|P}_i)}\times\underline{\mathrm O}(q_Q)
\]
(où, comme dessus, $\deleted{\unsure{q_Q}}Q=E_{\unsure{r|\nu}}^\perp/E_r$, $D=\{E_1\subset\ldots\subset E_r\}$). On constate que
% archive sheet number: 153
